%! Operations with Polynomials — Unit 4, Lesson 4.1 — Lesson Plan
\documentclass[10pt]{article}
\usepackage{algebra2-article}
\usepackage{algebra2-boxes}
\usepackage{graphicx}   % -article does NOT load graphicx; needed for thumbnails/screenshots
\graphicspath{{images/}}
\newcommand{\TallMath}[1]{$\displaystyle #1\rule[-1.4em]{0pt}{3.2em}$}
\newcommand{\CourseName}{Algebra 2: Shepherd}
\newcommand{\SchoolYear}{2026--2027}
\newcommand{\MeetingLength}{55 minutes}
\newcommand{\UnitNumberName}{Unit 4: Polynomial Functions \quad}
\newcommand{\LessonNumberName}{Lesson 4.1: Operations with Polynomials}

\begin{document}
\begin{center}
    {\Large\bfseries \CourseName: \SchoolYear \\
    {\small\bfseries \UnitNumberName \LessonNumberName}}
\end{center}

\begin{tcolorbox}[colback=forestbg, colframe=forest, boxrule=0.9pt, arc=2mm,
    left=2mm, right=2mm, top=1.5mm, bottom=1.5mm]
    \textbf{Primary Objective:} Students can \emph{add}, \emph{subtract}, and \emph{multiply}
    polynomial expressions in one and two variables and write the result in \emph{standard form} ---
    combining like terms, distributing a subtraction sign to \emph{every} term, multiplying term by
    term (monomial $\times$ polynomial, binomial $\times$ binomial, binomial $\times$ trinomial), and
    applying the \emph{special-product} identities $(a\pm b)^2$ and $(a+b)(a-b)$. They can then read
    the \emph{degree} and \emph{leading coefficient} of the answer --- degrees \emph{add} under
    multiplication, leading coefficients \emph{multiply} --- and use them to state the end behavior of
    the polynomial they just built. This turns Lesson~4.0's factored forms into standard forms and
    back.
    \\[2pt]
    \textbf{Standards (2023 VA SOL):} \textbf{A2.EO.3a} --- determine sums, differences, and products
    of polynomials in one \emph{and two} variables. \textbf{A2.EO.3d} appears in its \emph{forward}
    direction (demonstrating equality of polynomial expressions written in different forms, and the
    difference-of-squares and perfect-square-trinomial identities); the reverse direction ---
    factoring --- is Lesson~4.2. \textbf{A2.F.2g} is revisited whenever an expanded product's degree
    and leading coefficient are used to state end behavior.
\end{tcolorbox}

\renewcommand{\arraystretch}{1.4}
\begin{skillbox}[Priority Ideas \& Skills]{goldbox}
\begin{tabularx}{\linewidth}{@{}X|X@{}}
\textbf{Priority Ideas \& Skills} & \textbf{Key Understandings (the ``why'')} \\[2pt]
\hline
\vspace{-6pt}
\begin{itemize}[leftmargin=1.1em, nosep, after=\vspace{2pt}]
  \item Write a polynomial in \textbf{standard form} and name its \textbf{degree}, \textbf{leading coefficient}, and \textbf{number of terms}.
  \item \textbf{Add} and \textbf{subtract} polynomials by combining like terms --- distributing the subtraction to \emph{every} term.
  \item \textbf{Multiply} monomial $\times$ polynomial, binomial $\times$ binomial, and binomial $\times$ trinomial, in one and two variables.
  \item Apply the \textbf{special products} $(a+b)^2$, $(a-b)^2$, and $(a+b)(a-b)$.
  \item Predict a product's \textbf{degree} and \textbf{leading coefficient}, and state the resulting \textbf{end behavior}.
\end{itemize}
&
\vspace{-6pt}
\begin{itemize}[leftmargin=1.1em, nosep, after=\vspace{2pt}]
  \item Adding and subtracting only ever \emph{collects} terms; multiplying \emph{creates} new, higher-degree terms.
  \item A minus sign in front of a parenthesis is a factor of $-1$: it flips \emph{every} sign inside, not just the first.
  \item Multiplying is one rule --- \emph{every term times every term} --- whether it is called distributing, FOIL, or the box.
  \item Under multiplication degrees \textbf{add} and leading coefficients \textbf{multiply}, so the end behavior of a product is known before it is expanded.
  \item Factored form and standard form are the \emph{same function} in two costumes: one shows the zeros, the other shows the end behavior.
\end{itemize}
\end{tabularx}
\end{skillbox}

\begin{skillbox}[{Vocabulary, Concepts, \& Theorems}]{sky}
\renewcommand{\arraystretch}{1.3}
\begin{tabularx}{\linewidth}{@{}l X@{}}
\textbf{Standard form} & Terms written from the highest power down to the constant: $a_nx^n+\dots+a_1x+a_0$. \\
\textbf{Degree of a term} & The sum of the exponents on its variables --- so $-7x^2y^3$ has degree $5$. \\
\textbf{Degree / leading coefficient} & The largest term degree, and the coefficient of that term once in standard form. \\
\textbf{Monomial / binomial / trinomial} & A polynomial with one, two, or three terms. \\
\textbf{Like terms} & Terms with \emph{identical} variable parts (same variables, same exponents); only like terms combine. \\
\textbf{Distributive property} & $a(b+c)=ab+ac$ --- the single rule behind every polynomial product. \\
\textbf{Special products} & \TallMath{(a+b)^2=a^2+2ab+b^2,\quad (a-b)^2=a^2-2ab+b^2,\quad (a+b)(a-b)=a^2-b^2} \\
\textbf{Degree of a product} & $\deg(fg)=\deg f+\deg g$, and the leading coefficients multiply. \\
\textbf{Closure} & A sum, difference, or product of polynomials is always another polynomial. \\
\end{tabularx}
\end{skillbox}

\begin{fixedskillbox}[Activate Prior Knowledge \& Spiral Review]{forestbg}
    The Warm-Up rehearses exactly the four moves this lesson is built from: \textbf{(1)} the
    \emph{product of powers} rule $x^m\!\cdot\!x^n=x^{m+n}$ (multiply the coefficients, \emph{add} the
    exponents) --- the engine of every polynomial product; \textbf{(2)} \emph{combining like terms},
    which is all that addition and subtraction ever do; \textbf{(3)} \emph{distributing}, including
    the bare negative $-(4x^2-7x+3)$ that is this lesson's most common error; and \textbf{(4)} FOIL
    from Algebra~1, plus a one-line recall of \emph{degree} and \emph{leading coefficient} from
    Lesson~4.0. Debrief item~3 aloud --- ask what happened to the sign of \emph{each} of the three
    terms, not just the first.
\end{fixedskillbox}

\begin{skillbox}[Hook]{forestbg}
    Put both of Lesson~4.0's names for the anchor curve on the board:
    \[
      g(x) = (x-1)^2(x+2) \qquad\text{and}\qquad g(x)=x^3-3x+2 .
    \]
    \textbf{``We used both of these labels for the same graph yesterday. Are they really the same
    function --- or did we just take that on faith?''} Test it cheaply first: at $x=0$ both give $2$;
    at $x=2$ both give $4$; at $x=-2$ both give $0$. \textbf{``Agreeing at three points is good
    evidence --- but is it proof?''} (No: many different curves pass through three shared points.)
    Land the day's purpose: \emph{the only way to know two forms are the same function is to multiply
    the factored one out and see the standard one appear.} Note what each costume is good for --- the
    factored form hands you the \textbf{zeros and their multiplicities}, the standard form hands you
    the \textbf{degree and leading coefficient}, hence the \textbf{end behavior}. Today we learn to
    move from the first to the second.
\end{skillbox}

\begin{skillbox}[Lesson]{forestbg}
\begin{multicols}{2}
\textbf{1. Naming a polynomial.} Put it in \textbf{standard form} (powers high to low) first ---
every other name depends on it. Then read the \textbf{degree}, the \textbf{leading coefficient}, and
the \textbf{number of terms} (monomial / binomial / trinomial), and classify by degree (linear,
quadratic, cubic, quartic). For two-variable terms, the degree of a term is the \emph{sum} of its
exponents: $-7x^2y^3$ has degree $5$. Model on $-5x^2+7x^4-3+x \Rightarrow 7x^4-5x^2+x-3$.

\textbf{2. Adding \& subtracting: collecting, not creating.} \textbf{Like terms} must match in
\emph{both} variable and exponent; $3x^2$ and $3x^3$ never combine. Addition is pure collection. For
subtraction, rewrite the minus sign as a factor of $-1$ \emph{before} anything else --- \textbf{``the
minus sign belongs to every term in the parentheses.''} Show the same problem worked both ways
(distribute-then-collect vs.\ stacked in columns) and let students choose.

\textbf{Ask:} \textbf{``Can a sum of two degree-4 polynomials have degree less than 4?''} Yes ---
$(x^4+2x)-(x^4-5)$ leaves $2x+5$. The leading terms can cancel, so a sum or difference has degree
\emph{at most} the larger degree.

\columnbreak

\textbf{3. Multiplying: every term times every term.} One rule, three costumes. Monomial $\times$
polynomial is straight distribution ($-2x^2(3x^3-4x+6)$). Binomial $\times$ binomial is the same rule
with four products (FOIL is just a bookkeeping name). Binomial $\times$ trinomial has six products
--- draw the \textbf{box/area model} so nothing is dropped, then collect like terms down the
diagonals. Do at least one two-variable product, $(3x+2y)(x-4y)$, so the middle term $-10xy$ is seen
as a like-term collection.

\textbf{4. Degree \& leading coefficient of a product.} Before expanding $(2x^3-1)(-3x^4+x)$, ask
\textbf{``what will the highest-power term be?''} Only the two leading terms can produce it:
$2x^3\cdot(-3x^4)=-6x^7$. So \emph{degrees add, leading coefficients multiply} --- and the end
behavior from Lesson~4.0 (odd degree, negative lead $\Rightarrow$ left up, right down) is known
without doing the rest of the work.

\textbf{5. Special products.} Derive, don't memorize: $(a+b)^2=(a+b)(a+b)=a^2+2ab+b^2$. Flag the
year's most persistent error --- $(x+4)^2\ne x^2+16$ --- and kill it numerically at $x=1$
($25\ne17$). Then $(a-b)^2$ and the difference of squares $(a+b)(a-b)=a^2-b^2$, whose middle terms
cancel. Close by expanding the Hook's $(x-1)^2(x+2)$ to $x^3-3x+2$, and preview Lesson~4.2, where
these identities run \emph{backward}.
\end{multicols}
\end{skillbox}

\begin{skillbox}[Active Monitoring]{redbox}
Circulate during the notes practice and Tier work. Watch for and correct:
\begin{itemize}[leftmargin=1.2em, nosep]
  \item distributing a subtraction to only the \emph{first} term --- $(7x^2-4x+1)-(3x^2+6x-5)$ ending in $-5$ instead of $+6$;
  \item combining unlike terms ($3x^2+3x^3=6x^5$) --- the single most common ``simplification'' error;
  \item adding exponents when \emph{adding} terms, or multiplying them when \emph{multiplying} terms (the rule is: coefficients multiply, exponents add);
  \item $(a+b)^2$ written as $a^2+b^2$ --- send students to the $x=1$ numeric check rather than restating the rule;
  \item dropping a product in a binomial $\times$ trinomial --- require the box until all six cells appear;
  \item leaving the answer unordered --- standard form is part of the answer.
\end{itemize}
Cold-call prompts: ``What is the very first thing you do to a subtraction problem?'' ``What degree
will this product have --- before you expand it?'' ``Which two terms in the box are like terms, and
why?''
\end{skillbox}

\begin{skillbox}[Group Work \& Differentiation]{redbox}
\textbf{Building Polynomials} activity (tiered; mirrors the notes).
\begin{multicols}{3}
\textbf{Tier R --- Remediate.} One add, one subtract (with a ``rewrite with the signs changed'' line
built in), one monomial distribution, and one binomial product with the $2\times2$ box already
drawn. All answers in standard form.

\columnbreak
\textbf{Tier A --- Approaching.} A difference whose leading terms cancel (\emph{why} did the degree
drop?), a binomial $\times$ trinomial by box, two special products in \emph{two} variables, a
degree/leading-coefficient prediction with the end behavior it forces, and an error-analysis item on
$(x+4)^2=x^2+16$ requiring a numeric counterexample.

\columnbreak
\textbf{Tier E --- Extension.} Expand Lesson~4.0's open-box model $V(x)=x(10-2x)(8-2x)$ to
$4x^3-36x^2+80x$ and check it against that lesson's table; build a profit polynomial $P=R-C$ and
evaluate it in context; then justify why a \emph{product} of degree-3 and degree-4 polynomials always
has degree~7 while a \emph{sum} of two degree-4 polynomials need not have degree~4.
\end{multicols}
\end{skillbox}

\begin{skillbox}[Individual Work \& Assessment]{redbox}
\textbf{Exit Ticket} (independent, no notes): \textbf{(1)} a subtraction whose constant terms cancel
--- it is only correct if every sign was distributed; \textbf{(2)} a binomial $\times$ trinomial
product, followed by naming the result's degree and leading coefficient and \emph{stating the end
behavior} they force (the 4.0 link); and \textbf{(3)} an \textbf{SOL-style multiple-choice} item on
$(3x-4)^2$ whose distractors are exactly the three classic errors (dropped middle term, difference of
squares, and $2ab$ computed as $ab$). Collect and sort into ``got it / sign-distribution slip /
missing-middle-term slip / like-terms slip'' to decide whether Lesson~4.2 opens with a re-teach.
\end{skillbox}

\begin{skillbox}[Reinforcement \& Extension]{goldbox}
\textbf{Homework:} name a polynomial from standard form; add and subtract; multiply at all three
sizes (monomial, binomial $\times$ binomial, binomial $\times$ trinomial); apply the three special
products including a two-variable difference of squares; expand Lesson~4.0's exit-ticket function
$(x+1)^2(x-2)$ and confirm the end behavior matches the graph students already read; and one
justification item on why subtraction changes every sign. \textbf{Extension:} an area model --- a
$(2x+3)$ by $(x+5)$ garden with an $x$-by-$x$ pond removed --- requiring a product \emph{and} a
difference, then a numeric check. \textbf{Preview:} Lesson~4.2, \emph{Advanced Factoring} --- running
today's products backward, including the sum and difference of cubes.
\end{skillbox}
\end{document}
